% Part: first-order-logic
% Chapter: tableaux
% Section: derivations

\documentclass[../../../include/open-logic-section]{subfiles}

\begin{document}

\iftag{FOL}
      {\olfileid[fr]{fol}{tab}{der}}
      {\olfileid[fr]{pl}{tab}{der}}

\olsection{\usetoken{P}{tableau}}

\begin{explain}
Nous avons défini les hypothèses et donné les règles d'inférence.
Les !!{tableau}s sont engendrés inductivement à partir d'elles :
chaque !!{tableau} est soit une branche unique formée d'une ou plusieurs
hypothèses, soit le résultat de l'application d'une règle d'inférence
à une branche d'!!a{tableau}.
\end{explain}

\begin{defn}[!!^{tableau}]
!!^a{tableau} pour les hypothèses $\sFmla{S_1}{!A_1}$, \dots,
$\sFmla{S_n}{!A_n}$, où chaque $S_i$ vaut $\True$ ou~$\False$, est un
arbre fini de !!{signed formula}s qui vérifie les conditions suivantes :
\begin{enumerate}
\item Les $n$ !!{signed formula}s placées tout en haut de l'arbre sont
  $\sFmla{S_i}{!A_i}$, l'une au-dessous de l'autre.
\item Toute !!{signed formula} de l'arbre qui n'est pas une hypothèse
  résulte de l'application correcte d'une règle d'inférence à
  !!a{signed formula} située plus haut sur sa branche.
\end{enumerate}
Une branche d'!!a{tableau} est \emph{fermée} si et seulement si elle
contient à la fois $\sFmla{\True}{!A}$ et~$\sFmla{\False}{!A}$ ; elle
est \emph{ouverte} dans le cas contraire. !!^a{tableau} dont toutes les
branches sont fermées est un \emph{!!{tableau} fermé} pour son ensemble
d'hypothèses. Si !!a{tableau} n'est pas fermé, c'est-à-dire s'il contient
au moins une branche ouverte, il est \emph{ouvert}.
\end{defn}

\begin{ex}
  Tout ensemble d'hypothèses constitue à lui seul !!a{tableau}, mais
  celui-ci n'est généralement pas fermé. Il ne l'est manifestement que
  si les hypothèses contiennent déjà une paire de !!{signed formula}s
  $\sFmla{\True}{!A}$ et~$\sFmla{\False}{!A}$.

  À partir d'!!a{tableau}, ouvert ou fermé, on peut en obtenir un nouveau,
  plus grand, en appliquant l'une des règles d'inférence à
  !!a{signed formula} qui y figure et dont la formule est~$!A$. La règle ajoute une ou plusieurs
  !!{signed formula}s à l'extrémité de toute branche qui contient
  l'occurrence de~$!A$ à laquelle elle est appliquée.

  Considérons par exemple l'hypothèse $\sFmla{\True}{!A \land \lnot
    !A}$. Voici le !!{tableau} ouvert qui ne contient que cette hypothèse :
  \begin{center}
    \begin{tableau}{}
      [\sFmla{\True}{\formula{A} \land \lnot \formula{A}}, just=\TAss]
    \end{tableau}{}
  \end{center}
  On en obtient un nouveau !!{tableau} en appliquant la règle
  $\TRule{\True}{\land}$ à l'hypothèse. Cette règle permet d'ajouter
  deux lignes au !!{tableau}, $\sFmla{\True}{!A}$ et
  $\sFmla{\True}{\lnot !A}$ :
  \begin{center}
    \begin{tableau}{}
      [\sFmla{\True}{\formula{A} \land \lnot \formula{A}}, just=\TAss
        [\sFmla{\True}{\formula{A}}, just={\TRule{\True}{\land}[1]},
          [\sFmla{\True}{\lnot \formula{A}}, just={\TRule{\True}{\land}[1]}
          ]
        ]
      ]
    \end{tableau}{}
  \end{center}
  Lorsque nous écrivons des !!{tableau}s, nous indiquons à droite les
  règles appliquées. Ainsi, $\TRule{\True}{\land} 1$ signifie que la
  !!{signed formula} de cette ligne résulte de l'application de la règle
  $\TRule{\True}{\land}$ à la !!{signed formula} de la ligne~$1$.
  Ce nouveau !!{tableau} contient maintenant d'autres
  !!{signed formula}s, mais une seule d'entre elles,
  $\sFmla{\True}{\lnot !A}$, permet encore l'application d'une règle,
  en l'occurrence $\TRule{\True}{\lnot}$. On obtient ainsi le
  !!{tableau} fermé suivant :
  \begin{center}
    \begin{tableau}{}
      [\sFmla{\True}{\formula{A} \land \lnot \formula{A}}, just=\TAss
        [\sFmla{\True}{\formula{A}}, just={\TRule{\True}{\land}[1]}
          [\sFmla{\True}{\lnot \formula{A}}, just={\TRule{\True}{\land}[1]}
            [\sFmla{\False}{\formula{A}}, just={\TRule{\True}{\lnot}[3]}, close]
          ]
        ]
      ]
    \end{tableau}{}
  \end{center}
\end{ex}

\end{document}
